\lesson{1}{Sep 15 2021 Wed (13:40:07)}{Natural Rights}{Unit 1}

\subsection*{The Purpose of the Government}
\label{sub_sec:the_purpose_of_the_government}

The \textbf{Government} exists to provide order and safety to its citizens.
Without the \textbf{Government} and \textbf{Law}, chaos and anarchy would
reign.

The purposes of our \textbf{Government} is to keep the peace within its
territory, protect its citizens, provide important services, and maintain
public institution. The Government was created because collective
action, or groups of people working together, can accomplish these purposes
where individuals cannot. Laws and public policy help direct where and how to
build structures like airports. Things like speed limits, jail time, and other
policies are there to help protect the society. Institutions like schools and
hospitals help increase people's quality of life and improve society overall.
Government provides order to the creation and maintenance of institutions,
public policy, and infrastructure so that people can exercise their rights and
explore opportunities.

% subsection the_purpose_of_the_government (end)

\subsection*{How does the Government Carry out its Purpose}
\label{sub_sec:how_does_the_government_carry_out_its_purpose}

Governments pass laws about public policy to carry out their purposes in
specific situations. Laws are rules established by government or other sources
of authority to regulate people's conduct or activities. In the United States,
the U.S. Constitution is the basis for creating the following laws and policy.
This document outlines the structure and functions of the U.S. Government.
Americans follow the laws and policies within three levels of government:

\begin{itemize}
  \item Local
  \item State
  \item Federal
\end{itemize}

The Constitution guides us through the correct processes in the creation and
implementation of a law. Laws place limitations on citizen behavior. However,
most citizens agree with them because the law benefits the common good and
protects their most basic rights.

A citizen's duty is to abide by the laws. This helps protect order and the
rights of everyone. No one, not even the president, is above following the law.
Yet people do not always agree that a law is fair or appropriate. United States
citizens can challenge the laws through courts. Citizen consent to laws, and
upholding the rule of law affirms the legitimacy of the government.

% subsection how_does_the_government_carry_out_its_purpose (end)

\subsection*{What are Natural Rights}
\label{sub_sec:what_are_natural_rights}

The United States Constitution protects the natural rights of human beings. The
principle of natural rights has origins in ancient civilizations and religious
teachings, but the modern concept of it traces to the relatively recent
Enlightenment. John Locke described these rights as those of \textbf{"Life,
Liberty, and Estate"}. He used estate, or property, to explain that people have
the right to make their own living, to provide for their own needs. When the
Declaration of Independence was written, this last right was broadened to
"pursuit of happiness", as owning property was not guaranteed or always
necessary to provide for oneself.

Locke, and then the founders of the United States, believed God granted each
person these natural rights as individuals. Other philosophers agreed with the
idea of natural rights but said human beings had them because of their ability
to reason and act as rational beings. Either way, the idea of individual or
natural rights is a vital principle of the United States Government today.

The idea of individual, natural rights heavily influenced the founders of the
United States. It led them to include a list of specific protections for
rights, the \textbf{Bill of Rights}, as an addition section of the
Constitution. We can trace all of these rights to the ideas of \textbf{"Life,
Liberty, and Estate"}.

% subsection what_are_natural_rights (end)

\subsection*{What is Consent of the Governed}
\label{sub_sec:what_is_consent_of_the_governed}

A government is legitimate only if the people agree with its existence. Early
Americans declared their independence because they no longer consented to
British rule, as they felt the monarch did not respect or protect their rights
or their concerns. They believed that the only just government is one where the
people themselves, as free and equal citizens, determine what powers and
functions the government would have.

The related principle of social contract takes this a step further to say that
this consent places government under an obligation to fulfill its purpose by
protecting the people and their rights. If it does not, the people could
withdraw their consent, abolish the government, and form a new one. This quote
from the Declaration of Independence emphasizes this principle of citizen
consent and the social contract:

\begin{displayquote}
  \textit{That to secure these rights, governments are instituted among men, deriving
  their just powers from the consent of the governed. That whenever any form of
  government becomes destructive to these ends, it is the right of the people
  to alter or to abolish it, and to institute new government, laying its
  foundation on such principles.}
\end{displayquote}

% subsection what_is_consent_of_the_governed (end)

\subsection*{What is Representative Government}
\label{sub_sec:what_is_representative_government}

Democracy means \textbf{"Rules of the People"}. In history, two basic forms of
democracy existed -- direct and representative.

\begin{definition}[Direct Democracy]
  In Direct Democracy, every citizen would vote or consent to every decision.
  This is impractical in most societies.
\end{definition}

\begin{definition}[Representative Democracy]
  In Direct Representative, the people consent through their votes to giver
  certain people the power to make decisions on their behalf. Majority rule
  determines decisions, the law, or policy supported by the majority of people
  through their representatives. We also use the term \textbf{"Republic"} to
  describe a representative government. In the United States, most official
  positions carry a term of office where people have opportunity to choose
  someone else to represent them.
\end{definition}

% subsection what_is_representative_government (end)

\subsection*{What is Limited Government}
\label{sub_sec:what_is_limited_government}

The founders of the United States wanted to prevent a government growing too
large, where it would become unrestrained and infringe on people's rights. They
believed it should be as small as possible and its powers limited to certain
functions, any protection, or service that affected the nation as a whole.

They clearly outlined the structure, functions, and powers of the United States
through the Constitution. While the government has expanded or contracted its
powers over time, the people and representatives have had a voice in most
changes, ensuring that the government continues to only use powers which the
people consent are necessary.

% subsection what_is_limited_government (end)

\subsection*{What is Rule of Law}
\label{sub_sec:what_is_rule_of_law}

This fundamental principle means that no one is above or below the law. That
means that if a law is broken, no matter by whom, certain procedures will be
followed. Even the president and other federal officials must abide by the
laws. They pay taxes and obey speed limits like everyone else. Laws exist to
protect safety and order, and all citizens must obey them. It's their duty.

% subsection what_is_rule_of_law (end)

\newpage
